\documentclass[a4paper]{article}     
\usepackage{amsfonts}
\usepackage{amsmath}
\usepackage{amssymb}
\usepackage{graphicx}
\usepackage{color}

\newcommand{\Q}{\mathbb{Q}}
\newcommand{\R}{\mathbb{R}}
\newcommand{\llim}{\lim\limits}
\newcommand{\dint}{\displaystyle\int}

\begin{document}

\noindent \large Auteur: Hina Rana \\
\noindent \large Studierichting: Informatica\\
\noindent \large Oefeningenreeks 6B nr. 17.1 \\

\medskip

\normalsize

$ n(n+1) \rightarrow S_k  = k(k+1) = k^2 + k $ \\

$ s_n = \displaystyle\sum_{k=1}^{n} k^2 + k = \displaystyle\sum_{k=1}^{n} k^2 + \sum_{k=1}^{n} k $ \\

$ s_n = S_2(n) + S_1(n) $ \\

$ s_n = \left(\dfrac{1}{3} n^3 + \dfrac{1}{2} n^2 + \dfrac{1}{6} n \right) + \left(\dfrac{1}{2} n^2 + \dfrac{1}{2} n \right) $ \\

$ = \dfrac{1}{3} n^3 + 1n^2 + \dfrac{2}{3} n $ \\

$ s_n = \dfrac{n^3 + 3n^2 + 2n}{3} $ \\

$ s_n = \dfrac{n(n^2 + 3n + 2)}{3} $ \\

$ s_n = \dfrac{n(n+1)(n+2)}{3} $ \\

Dit is een natuurlijk getal want een product van 3 opeenvolgende getallen is deelbaar door 3.

\begin{thebibliography}{999}
\end{thebibliography}
\end{document} 
